\section{Alternativas de diseño propuestas en el proceso}
\label{sec:alternativas}

Esta sección documenta las decisiones de diseño más relevantes tomadas
durante la conceptualización del proyecto. Para cada una de las cinco
dimensiones críticas se presentan las alternativas evaluadas, se
comparan bajo los mismos criterios (usabilidad, costo de
implementación, accesibilidad y riesgo técnico) y se justifica la
opción seleccionada con trazabilidad explícita a los requerimientos
de la sección~\ref{sec:requerimientos} y a los casos de uso de la
sección~\ref{sec:casos}.

\subsection{Dimensión 1: representación de las métricas durante la entrevista}

El reto de diseño consistió en comunicar al candidato información
continua sobre su desempeño (fluidez, contacto visual, postura) sin
desviar su atención del diálogo con el entrevistador.

\subsubsection*{Alternativa 1A — Paneles laterales numéricos}
Disposición convencional con el video del candidato al centro y
paneles laterales con barras y cifras (``contacto visual: 62\,\%'',
``muletillas: 3/min''), análoga a la usada por Yoodli o Big Interview.

\begin{table}[ht]
\centering\footnotesize
\renewcommand{\arraystretch}{1.2}
\begin{tabularx}{\linewidth}{@{}p{3.6cm} Y@{}}
\toprule
\textbf{Criterio} & \textbf{Evaluación} \\
\midrule
Usabilidad & Diseño familiar y de bajo costo cognitivo inicial, pero
divide la atención entre entrevistador y paneles, interrumpe el ritmo
conversacional y puede aumentar la ansiedad por exposición continua a
cifras. \\
Costo de implementación & Bajo; patrón de UI estándar. \\
Accesibilidad & Compatible con lectores de pantalla mediante etiquetas
ARIA, pero la carga visual adicional perjudica al perfil con
neurodivergencia (\sk{7}). \\
Riesgo técnico & Muy bajo. \\
\bottomrule
\end{tabularx}
\end{table}

\noindent\textbf{Decisión:} descartada. No cumple el módulo de
interfaz gráfica no convencional exigido por el enunciado y rompe la
inmersión propia de una entrevista.

\subsubsection*{Alternativa 1B — Indicadores radiales superpuestos}
Círculos de progreso superpuestos al avatar que cambian de tamaño o
color según las métricas. Punto intermedio entre paneles y aura.

\begin{table}[ht]
\centering\footnotesize
\renewcommand{\arraystretch}{1.2}
\begin{tabularx}{\linewidth}{@{}p{3.6cm} Y@{}}
\toprule
\textbf{Criterio} & \textbf{Evaluación} \\
\midrule
Usabilidad & Menos intrusivos que los paneles, pero requieren mirada
directa para ser leídos; no se perciben de forma periférica. \\
Costo de implementación & Medio; viable con SVG animado o
transiciones CSS. \\
Accesibilidad & Si la información se transmite únicamente por color,
incumple WCAG 2.2 criterio 1.4.1, por lo que requiere un canal
adicional. \\
Riesgo técnico & Bajo. \\
\bottomrule
\end{tabularx}
\end{table}

\noindent\textbf{Decisión:} descartada. La lectura no es periférica y
la dependencia exclusiva del color genera barreras de accesibilidad.

\subsubsection*{Alternativa 1C — Aura luminosa procedural alrededor del avatar (elegida)}
Halo visual generado con WebGL/Three.js que rodea al avatar y modula
sus parámetros en función de las métricas: color (valencia global),
intensidad (fluidez verbal), pulsación (cadencia) y asimetría
espacial (equilibrio entre canal verbal y no verbal). Durante la
entrevista no hay paneles ni números: la información reside en la
periferia visual.

\begin{table}[ht]
\centering\footnotesize
\renewcommand{\arraystretch}{1.2}
\begin{tabularx}{\linewidth}{@{}p{3.6cm} Y@{}}
\toprule
\textbf{Criterio} & \textbf{Evaluación} \\
\midrule
Usabilidad & Permite mantener el contacto visual con el avatar y
recibir simultáneamente información del desempeño. El tour inicial y
los \textit{tooltips} contextuales (\rf{15}) cubren la curva de
aprendizaje del significado de cada parámetro. \\
Costo de implementación & Alto; requiere Three.js y \textit{shaders}
GLSL. La degradación automática a una representación bidimensional
simplificada cuando se detecta GPU insuficiente preserva el
rendimiento (\rnf{03}). \\
Accesibilidad & El canal visual del aura cuenta con un equivalente
sonoro mediante \textit{earcons} (\rf{17}) para usuarios con
discapacidad visual (\sk{4}); el aura nunca es el único canal. \\
Riesgo técnico & Medio; mitigado mediante la ruta de degradación 2D y
la verificación de la latencia del aura ($\le 500$\,ms p95,
\rnf{01}). \\
\bottomrule
\end{tabularx}
\end{table}

\noindent\textbf{Decisión:} \textbf{opción elegida}. Es la única
alternativa que entrega \textit{feedback} continuo sin interrumpir el
diálogo y la única que materializa el módulo de interfaz no
convencional. Cubre \rf{05}, \rf{15}, \rnf{01}, \rnf{03}.

\subsection{Dimensión 2: representación visual del entrevistador}

Se evaluó qué forma debía adoptar el entrevistador en pantalla para
ser percibido como interlocutor sin generar incomodidad.

\subsubsection*{Alternativa 2A — Solo voz, sin figura visual}
El entrevistador existe únicamente como canal de audio. En pantalla
permanecen el aura, los controles de la sesión y la cámara del
candidato.

\begin{table}[ht]
\centering\footnotesize
\renewcommand{\arraystretch}{1.2}
\begin{tabularx}{\linewidth}{@{}p{3.6cm} Y@{}}
\toprule
\textbf{Criterio} & \textbf{Evaluación} \\
\midrule
Usabilidad & La ausencia de un punto visual de referencia hace
artificial el intercambio respecto a una entrevista presencial y
debilita la metáfora del \textit{Warachikuy}. \\
Costo de implementación & Muy bajo; no hay avatar que renderizar. \\
Accesibilidad & Favorable para \sk{4}, pero perjudica a \sk{5} por la
ausencia de señales visuales del cambio de turno. \\
Riesgo técnico & Muy bajo. \\
\bottomrule
\end{tabularx}
\end{table}

\noindent\textbf{Decisión:} descartada como diseño principal. Se
mantiene como modo de degradación cuando el hardware no soporta el
renderizado 3D.

\subsubsection*{Alternativa 2B — Avatar 2D estilo \textit{cartoon}}
Ilustración animada bidimensional con sincronía básica de boca y
ojos, similar a aplicaciones formativas masivas.

\begin{table}[ht]
\centering\footnotesize
\renewcommand{\arraystretch}{1.2}
\begin{tabularx}{\linewidth}{@{}p{3.6cm} Y@{}}
\toprule
\textbf{Criterio} & \textbf{Evaluación} \\
\midrule
Usabilidad & Formato familiar, pero el registro estilístico choca con
la formalidad esperada de una entrevista laboral y proyecta una
imagen poco profesional. \\
Costo de implementación & Medio; viable con Lottie o Rive. \\
Accesibilidad & Más liviano que el 3D y compatible con hardware
básico sin necesidad de degradación. \\
Riesgo técnico & Bajo. \\
\bottomrule
\end{tabularx}
\end{table}

\noindent\textbf{Decisión:} descartada. El registro estilístico no es
coherente con la simulación de un proceso de selección real.

\subsubsection*{Alternativa 2C — Avatar 3D estilizado, no hiperrealista (elegida)}
Personaje 3D con geometría simplificada y estética no fotorrealista
(\textit{cel-shading} o \textit{low-poly}) renderizado con Three.js.
La estilización evita el \textit{uncanny valley} y al mismo tiempo
preserva la presencia visual de un interlocutor; el aura procedural
se integra como una extensión natural del avatar.

\begin{table}[ht]
\centering\footnotesize
\renewcommand{\arraystretch}{1.2}
\begin{tabularx}{\linewidth}{@{}p{3.6cm} Y@{}}
\toprule
\textbf{Criterio} & \textbf{Evaluación} \\
\midrule
Usabilidad & La abstracción visual es suficiente para que el usuario
trate al avatar como un interlocutor sin la incomodidad asociada a
representaciones casi-realistas. Encaja con la metáfora del
\textit{Warachikuy} adoptada por el proyecto. \\
Costo de implementación & Alto; mitigado mediante el uso de
\textit{assets} libres (ReadyPlayerMe, Mixamo) y animaciones
genéricas para el habla. \\
Accesibilidad & Cuando el hardware no soporta el 3D, se activa
automáticamente la versión 2D o solo-voz (\rnf{03}). \\
Riesgo técnico & Medio; el principal foco es la sincronía
voz-animación, simplificable a una animación genérica de turno. \\
\bottomrule
\end{tabularx}
\end{table}

\noindent\textbf{Decisión:} \textbf{opción elegida}. Es coherente con
la metáfora del proyecto, permite una integración orgánica con el
aura y dispone de una ruta de degradación controlada. Cubre \rf{05}
y \rnf{03}.

\subsection{Dimensión 3: tecnología de renderizado del aura}

Una vez definido el aura como canal visual, se evaluó la tecnología
adecuada para producirla, considerando equipos sin GPU dedicada.

\subsubsection*{Alternativa 3A — Canvas 2D nativo del navegador}
El aura se dibuja como formas bidimensionales (círculos, ondas)
animadas sobre un \textit{canvas} plano superpuesto al avatar.

\begin{table}[ht]
\centering\footnotesize
\renewcommand{\arraystretch}{1.2}
\begin{tabularx}{\linewidth}{@{}p{3.6cm} Y@{}}
\toprule
\textbf{Criterio} & \textbf{Evaluación} \\
\midrule
Usabilidad & La planitud impide transmitir la noción de halo
envolvente; la composición resulta en una capa adherida en lugar de
una integración orgánica. \\
Costo de implementación & Bajo; API estándar bien documentada. \\
Riesgo técnico & Muy bajo, pero la discontinuidad visual entre el
plano del aura y el entorno tridimensional del avatar deteriora la
coherencia. \\
\bottomrule
\end{tabularx}
\end{table}

\noindent\textbf{Decisión:} descartada por incompatibilidad visual
con el avatar 3D.

\subsubsection*{Alternativa 3B — SVG animado con GSAP}
El aura se compone de figuras vectoriales SVG cuyos atributos se
animan mediante CSS o GSAP.

\begin{table}[ht]
\centering\footnotesize
\renewcommand{\arraystretch}{1.2}
\begin{tabularx}{\linewidth}{@{}p{3.6cm} Y@{}}
\toprule
\textbf{Criterio} & \textbf{Evaluación} \\
\midrule
Usabilidad & Los efectos son más expresivos que en Canvas 2D, pero
persiste el problema de integración con el entorno 3D; resulta
adecuado como modo de respaldo. \\
Costo de implementación & Bajo-medio; GSAP cuenta con
documentación y ejemplos abundantes. \\
Riesgo técnico & Bajo; con muchos \textit{paths} SVG simultáneos se
observan caídas de rendimiento. \\
\bottomrule
\end{tabularx}
\end{table}

\noindent\textbf{Decisión:} descartada como opción principal. Se
adopta como \textit{fallback} 2D cuando se detecta ausencia de GPU,
sustituyendo al aura WebGL en ese caso.

\subsubsection*{Alternativa 3C — WebGL/Three.js con \textit{shaders} (elegida)}
El aura se implementa como una esfera deformada en tiempo real
mediante \textit{shaders} GLSL, integrada en la misma escena 3D del
avatar. Los parámetros de deformación, color y ritmo se actualizan
desde el bucle de métricas multimodales.

\begin{table}[ht]
\centering\footnotesize
\renewcommand{\arraystretch}{1.2}
\begin{tabularx}{\linewidth}{@{}p{3.6cm} Y@{}}
\toprule
\textbf{Criterio} & \textbf{Evaluación} \\
\midrule
Usabilidad & El aura queda integrada de forma orgánica al avatar; los
cambios de color e intensidad se perciben naturalmente en la
periferia visual. \\
Costo de implementación & Alto; requiere conocimientos de GLSL y
Three.js. Mitigado partiendo de \textit{shaders} de plasma de
referencia disponibles en repositorios abiertos. \\
Accesibilidad & El canal visual cuenta con un equivalente sonoro
permanente para \sk{4} (\rf{17}). \\
Riesgo técnico & Medio; cubierto por la degradación automática a la
alternativa 3B cuando el hardware lo requiere. \\
\bottomrule
\end{tabularx}
\end{table}

\noindent\textbf{Decisión:} \textbf{opción elegida}. Es la única que
produce un halo orgánico integrado al avatar 3D, con una ruta de
degradación que controla el riesgo de hardware. Cubre \rf{05},
\rnf{01} y \rnf{03}.

\subsection{Dimensión 4: modalidad del \textit{peer mock}}

Se evaluaron tres modalidades para el módulo de práctica entre pares,
en función de su sincronía y de la riqueza de \textit{feedback}
disponible para el observador.

\subsubsection*{Alternativa 4A — Asincrónica con grabación y comentarios diferidos}
El candidato graba la sesión y la comparte con un par que comenta en
\textit{timestamps} específicos sin necesidad de coincidir
temporalmente.

\begin{table}[ht]
\centering\footnotesize
\renewcommand{\arraystretch}{1.2}
\begin{tabularx}{\linewidth}{@{}p{3.6cm} Y@{}}
\toprule
\textbf{Criterio} & \textbf{Evaluación} \\
\midrule
Usabilidad & Flexible, no requiere coordinación de horarios; pierde
la presión propia de hablar en vivo, que es parte central del
entrenamiento. \\
Costo de implementación & Medio; requiere almacenamiento S3 y
sistema de comentarios anclados, sin WebRTC. \\
Accesibilidad & Beneficiosa para \sk{6} y \sk{7} cuando requieren
mayor tiempo de procesamiento. \\
Riesgo técnico & Bajo; sin la complejidad de la conectividad P2P. \\
\bottomrule
\end{tabularx}
\end{table}

\noindent\textbf{Decisión:} no elegida como modalidad principal; se
incorpora como funcionalidad complementaria mediante los comentarios
anclados a \textit{timestamps} (\rf{11}) dentro del flujo síncrono.

\subsubsection*{Alternativa 4B — Videollamada externa sin métricas compartidas}
Los usuarios se conectan por una plataforma de terceros (Zoom o Meet)
mientras el candidato realiza la entrevista con el LLM; el observador
asiste sin acceso a las métricas en vivo.

\begin{table}[ht]
\centering\footnotesize
\renewcommand{\arraystretch}{1.2}
\begin{tabularx}{\linewidth}{@{}p{3.6cm} Y@{}}
\toprule
\textbf{Criterio} & \textbf{Evaluación} \\
\midrule
Usabilidad & Familiaridad con la plataforma, pero el observador no
accede al aura ni puede anclar comentarios; se diluye el valor
diferencial del \textit{peer mock}. \\
Costo de implementación & Bajo en integración; alto en dependencia
externa. \\
Accesibilidad & Sujeta a la accesibilidad de la plataforma de
terceros, fuera del control del proyecto. \\
Riesgo técnico & Bajo en lo técnico, alto en gobernanza y pérdida de
funcionalidad. \\
\bottomrule
\end{tabularx}
\end{table}

\noindent\textbf{Decisión:} descartada. La ausencia de métricas
compartidas vacía de contenido al módulo.

\subsubsection*{Alternativa 4C — Síncrona con WebRTC y panel del observador en vivo (elegida)}
Conexión \textit{peer-to-peer} en tiempo real vía WebRTC con
señalización por servidor propio. Un usuario asume el rol candidato y
otro el de observador con un panel lateral que muestra las métricas
en vivo y permite anclar comentarios; al cierre, los pares pueden
intercambiar roles.

\begin{table}[ht]
\centering\footnotesize
\renewcommand{\arraystretch}{1.2}
\begin{tabularx}{\linewidth}{@{}p{3.6cm} Y@{}}
\toprule
\textbf{Criterio} & \textbf{Evaluación} \\
\midrule
Usabilidad & Reproduce la presión real de la entrevista; el
intercambio de roles desarrolla la perspectiva del entrevistador. \\
Costo de implementación & Alto; requiere servidor de señalización,
servidor TURN como respaldo, sincronía de métricas y panel del
observador. \\
Accesibilidad & El canal de chat textual paralelo garantiza la
participación de \sk{5}; el consentimiento de grabación es
obligatorio antes del inicio (\rf{21}). \\
Riesgo técnico & Alto en conectividad P2P; mitigado con TURN y un
modo asincrónico de respaldo cuando la conexión falla
(\cu{03}, flujo alternativo 2a). \\
\bottomrule
\end{tabularx}
\end{table}

\noindent\textbf{Decisión:} \textbf{opción elegida}. La práctica
sincrónica con observador instrumentado es el diferencial del
proyecto frente al estado del arte y resulta coherente con el
objetivo de entrenar bajo presión real. Cubre \rf{10}, \rf{11},
\rf{12} y el caso de uso \cu{03}.

\subsection{Dimensión 5: adaptación de la dificultad de las preguntas}

El módulo de personalización continua exige que la dificultad de las
preguntas se ajuste al nivel del candidato. Se evaluaron tres
estrategias.

\subsubsection*{Alternativa 5A — Nivel fijo declarado por el usuario}
El usuario selecciona su nivel (\textit{junior}, \textit{mid},
\textit{senior}) al inicio y permanece constante en todas las
sesiones.

\begin{table}[ht]
\centering\footnotesize
\renewcommand{\arraystretch}{1.2}
\begin{tabularx}{\linewidth}{@{}p{3.6cm} Y@{}}
\toprule
\textbf{Criterio} & \textbf{Evaluación} \\
\midrule
Usabilidad & Predecible para el usuario; ignora la mejora longitudinal
y entrega preguntas de dificultad estática. \\
Costo de implementación & Muy bajo. \\
Riesgo técnico & Muy bajo. \\
\bottomrule
\end{tabularx}
\end{table}

\noindent\textbf{Decisión:} descartada. Contradice el módulo de
personalización continua exigido por el enunciado.

\subsubsection*{Alternativa 5B — Reglas heurísticas escritas por el equipo}
Conjunto de reglas fijas (por ejemplo, ``si el puntaje de las últimas
tres respuestas supera 80, subir nivel'') que ajustan la dificultad
de forma determinista.

\begin{table}[ht]
\centering\footnotesize
\renewcommand{\arraystretch}{1.2}
\begin{tabularx}{\linewidth}{@{}p{3.6cm} Y@{}}
\toprule
\textbf{Criterio} & \textbf{Evaluación} \\
\midrule
Usabilidad & Comportamiento predecible si las reglas son
transparentes; las reglas globales no distinguen entre competencias,
de modo que un candidato con desempeño dispar entre comunicación y
contenido técnico recibe una calibración única. \\
Costo de implementación & Bajo-medio; lógica de negocio sin modelos
estadísticos. \\
Riesgo técnico & Bajo; la mantenibilidad se degrada al crecer el
catálogo de industrias y roles. \\
\bottomrule
\end{tabularx}
\end{table}

\noindent\textbf{Decisión:} descartada como modelo principal. Se
emplea durante la primera sesión, mientras el sistema acumula
historial suficiente para activar el modelo más complejo.

\subsubsection*{Alternativa 5C — Adaptación con Teoría de Respuesta al Ítem (IRT) (elegida)}
Modelo estadístico basado en IRT que estima la habilidad del candidato
por competencia (comunicación verbal, contenido técnico, manejo del
estrés, lenguaje corporal) de forma independiente, se actualiza tras
cada respuesta y selecciona la siguiente pregunta según el nivel
estimado.

\begin{table}[ht]
\centering\footnotesize
\renewcommand{\arraystretch}{1.2}
\begin{tabularx}{\linewidth}{@{}p{3.6cm} Y@{}}
\toprule
\textbf{Criterio} & \textbf{Evaluación} \\
\midrule
Usabilidad & Adaptación precisa por competencia: el sistema puede
formular preguntas de mayor dificultad técnica y de menor dificultad
en manejo del estrés en una misma sesión, según el desempeño. \\
Costo de implementación & Alto; requiere parametrizar el banco de
preguntas con dificultad y discriminación, e implementar el estimador
de habilidad. Se adopta una versión inicial con modelo 1PL
(Rasch), evolutiva hacia 2PL/3PL. \\
Accesibilidad & El modelo opera sobre métricas derivadas, no sobre
video crudo, por lo que no penaliza a \sk{4}--\sk{7} cuando utilizan
las métricas alternativas previstas para sus rutas accesibles. \\
Riesgo técnico & Medio-alto al inicio por escasez de datos; mitigado
arrancando con la alternativa 5B y activando el IRT cuando se
alcanza un volumen mínimo de sesiones registradas. \\
\bottomrule
\end{tabularx}
\end{table}

\noindent\textbf{Decisión:} \textbf{opción elegida}. Es la única que
realiza personalización por competencia y mejora con el uso. La
implementación adopta el modelo 1PL como punto de partida. Cubre
\rf{09} y el módulo opcional de personalización continua.

\subsection{Síntesis comparativa}

La tabla \ref{tab:alternativas-resumen} consolida las cinco
dimensiones, las opciones evaluadas y la decisión adoptada, con la
trazabilidad agregada a requerimientos.

\begin{table}[ht]
\centering\footnotesize
\renewcommand{\arraystretch}{1.25}
\begin{tabularx}{\linewidth}{@{}p{3.4cm} Y p{3.0cm} p{3.0cm}@{}}
\toprule
\textbf{Dimensión} & \textbf{Opciones evaluadas} &
\textbf{Decisión} & \textbf{Trazabilidad} \\
\midrule
1. Métricas en sesión & Paneles laterales · Indicadores radiales ·
Aura procedural & Aura procedural & \rf{05}, \rf{15}, \rnf{01},
\rnf{03} \\
\addlinespace
2. Representación del entrevistador & Solo voz · Avatar 2D · Avatar 3D
estilizado & Avatar 3D estilizado & \rf{05}, \rnf{03} \\
\addlinespace
3. Tecnología del aura & Canvas 2D · SVG con GSAP · WebGL/Three.js &
WebGL/Three.js (\textit{fallback} SVG) & \rf{05}, \rnf{01}, \rnf{03} \\
\addlinespace
4. \textit{Peer mock} & Asincrónico · Videollamada externa · WebRTC
con panel del observador & WebRTC con observador instrumentado &
\rf{10}--\rf{12}, \cu{03} \\
\addlinespace
5. Adaptación de dificultad & Nivel fijo · Reglas heurísticas · IRT
por competencia & IRT por competencia (1PL inicial) & \rf{09} \\
\bottomrule
\end{tabularx}
\caption{Síntesis de alternativas evaluadas y decisiones adoptadas
por dimensión.}
\label{tab:alternativas-resumen}
\end{table}
